\printedpage{4}
% head: A. Grothendieck : Dérivateurs, Ch. IX, Édition des Universités de Montpellier II et Paris VII
% foot: Version 12 mars 2002

we deduce from this
\[
\left\{
\begin{array}{lcl}
\pi_0(P) & \xrightarrow{c_u^{\mathrm{l}}} & \mathrm{Hom}_{\mathcal{H}}(Q,R)\\
\pi_0(Q) & \xrightarrow{c_u^{\mathrm{r}}} & \mathrm{Hom}_{\mathcal{H}}(P,R)
\end{array}
\right.
\]
as follows:
\[
\begin{array}{rlcl}
\text{If } \alpha\in\pi_0(P),\ \text{i.e. } 1\xrightarrow{\ \alpha\ }P,\ \text{then} & \multicolumn{3}{l}{\alpha*\mathrm{id}_Q : Q\longrightarrow P*Q}\\
\text{and} & c^{\mathrm{l}}_u(\alpha) & = & u\circ(\alpha*\mathrm{id}_Q)\\
\text{and dually} & c^{\mathrm{r}}_u(\beta) & = & u\circ(\mathrm{id}_P*\beta).
\end{array}
\]
If we have three objects $P_0$, $P_1$, $P_2$ in $\mathcal{H}$, pairings
\[
\left\{
\begin{array}{rcl}
P_0*P_1 & \xrightarrow{u_1} & Q_0\\
P_1*P_2 & \xrightarrow{u_2} & Q_1\\
P_0*Q_1 & \xrightarrow{u_3} & R\\
Q_0*P_2 & \xrightarrow{u_4} & R
\end{array}
\right.
\]
\[
\begin{tikzcd}
a \arrow[r, "P_0"'] \arrow[rr, bend left=20, "Q_0"] \arrow[rrr, bend left=35, "R"] & b \arrow[r, "P_1"'] \arrow[rr, bend right=20, "Q_1"'] & c \arrow[r, "P_2"'] & d
\end{tikzcd}
\]

\msp{page 6}
with an associativity relation, saying that
\[
\begin{tikzcd}
P_0*P_1*P_2 \arrow[r] \arrow[d] & Q_0*P_2 \arrow[d]\\
P_0*Q_1 \arrow[r] & R
\end{tikzcd}
\]
is commutative, then we deduce from this two transitivity relations and one commutation relation, for the operations associated with the elements of $\pi_0(P_0)$, $\pi_0(P_1)$, $\pi_0(P_2)$, which we will not spell out, except for one:
\[
\begin{array}{rrrcl}
\multicolumn{5}{l}{\text{Let } \beta\in P_1,\ \gamma\in P_2,}\\
\text{whence by } \beta : & c_{u_1}(\beta) : & P_0 & \longrightarrow & Q_0,\\
\text{and by } \gamma : & c_{u_4}(\gamma) : & Q_0 & \longrightarrow & R,\\
\text{and by } \beta *_{u_2}\gamma\in\pi_0(Q_1) : & c_{u_3}(\beta\,\texttt{[}{*_{u_2}}\texttt{]}\gamma) : & P_0 & \longrightarrow & R,
\end{array}
\]

\printedpage{5}
% head: A. Grothendieck : Dérivateurs, Ch. IX, Édition des Universités de Montpellier II et Paris VII
% foot: Version 12 mars 2002
\texttt{[rather} $\beta\in\pi_0(P_1)$\texttt{,} $\gamma\in\pi_0(P_2)$\texttt{]}, well, we have
\[
\boxed{\;c(\beta *_{u_2}\gamma) \;=\; c(\gamma)c(\beta).\;}
\]
We still have to define $\beta *_{u_2}\gamma$, via the pairing $P_1*P_2\xrightarrow{u_2}Q_1$, we have in fact
\[
\begin{array}{ccc}
\pi_0(P_1)\times\pi_0(P_2) & \longrightarrow & \pi_0(P_1*P_2)\\
(\beta\,,\,\gamma) & \longmapsto & \beta*\gamma,
\end{array}
\]
and we compose with
\[
\pi_0(u_2) : P_1*P_2 \longrightarrow Q_1
\]
\texttt{[rather} $\pi_0(u_2) : \pi_0(P_1*P_2) \longrightarrow \pi_0(Q_1)$\texttt{]}.

\msp{page 7}
We give ourselves a set $\mathcal{M}_0$, and
\[
\left\{
\begin{array}{rclccrcl}
\mathcal{M}_0\times\mathcal{M}_0 & \xrightarrow{\mathrm{Ch}} & \mathcal{H} & \text{denoted} & & (a,b) & \longmapsto & \mathrm{Ch}(a,b)\ \text{or}\ \mathrm{Hom\,ext}_{\mathcal{H}}(a,b)\\
\mathcal{M}_0\times\mathcal{M}_0\times\mathcal{M}_0 & \xrightarrow{\text{comp.}} & \mathrm{Fl}\,\mathcal{H} & & & (a,b,c) & \longmapsto & \bigl(\mathrm{Ch}(a,b)\times\mathrm{Ch}(b,c)\longrightarrow\mathrm{Ch}(a,c)\bigr)
\end{array}
\right.
\]
with two axioms

\begin{enumerate}
\item[a)] Associativity for the compositions, for $(a,b,c,d)\in\mathcal{M}_0^4$
\[
\begin{tikzcd}
\mathrm{Ch}(a,b)\times\mathrm{Ch}(b,c)\times\mathrm{Ch}(c,d) \arrow[r] \arrow[d] & \mathrm{Ch}(a,c)\times\mathrm{Ch}(c,d) \arrow[d]\\
\mathrm{Ch}(a,b)\times\mathrm{Ch}(b,d) \arrow[r] & \mathrm{Ch}(a,d).
\end{tikzcd}
\]
\item[b)] Existence of two-sided unit objects
\[
1_a\in\pi_0(\mathrm{Ch}(a,a)) \qquad \text{for every } a\in\mathcal{M}_0
\]
\end{enumerate}

Then the $\pi_0(\mathrm{Ch}(a,b))$ and their composition define on $\mathcal{M}_0$ as set of objects a \textit{category} structure $\MOP$, having the $1_a$ as identities, and we find that
\[
(a,b)\longmapsto\mathrm{Ch}(a,b)
\]
can be interpreted in terms of a canonical \textit{bifunctor}
\[
\mathcal{M}^o_{0\,\Pi}\times\MOP \longrightarrow \mathcal{H},
\]
whose composite with
\[
\mathcal{H}\xrightarrow{\pi_0}\mathrm{Ens}
\]
is none other than the bifunctor $(a,b)\longrightarrow\mathrm{Hom}_{\MOP}(a,b)$,
\printedpage{6}
% head: A. Grothendieck : Dérivateurs, Ch. IX, Édition des Universités de Montpellier II et Paris VII
% foot: Version 12 mars 2002
\msp{page 8}
\[
\mathrm{Hom}_{\MOP}(a,b) \;\simeq\; \pi_0(\mathrm{Ch}(a,b)).
\]
This relation is the \textit{definition} of $\underline{\mathrm{Hom}}_{\MOP}$, but here the relation is understood as a % Reading resolved against PDF leaf 18: a single underline; see corrections.csv.
\textit{functorial} bijection in $a$, $b$.

We find in particular that if
\[
u : a'\longrightarrow a, \qquad v : b\longrightarrow b'
\]
are \textit{invertible} arrows in $\MOP$, then
\[
\mathrm{Ch}(a,b)\longrightarrow\mathrm{Ch}(a',b')
\]
corresponding in $\mathcal{H}$ is \texttt{[an]} \textit{isomorphism.}

We will apply all this mainly to the two situations deduced from a situation $(\mathcal{M},\Sigma)$, with
\[
\begin{array}{rlcllcl}
1^o) & \mathcal{H} & = & \mathrm{Hot}, & \mathrm{Ch}(a,b) & = & \mathrm{Ch}^{\Sigma}_{\infty}(\mathcal{M};a,b)\\
2^o) & \mathcal{H} & = & \mathrm{Ind\,Cat\,hot}, & \mathrm{Ch}(a,b) & = & \mathrm{CH}^{\Sigma}_{\infty}(\mathcal{M};a,b)\ \text{(\footnotemark[2])}.
\end{array}
\]
\footnotetext[2]{$*$ in Hot and Ind\,Cat\,hot is the cartesian product, nothing more.}
In both cases, we thus have a functor on $(\mathcal{M}\Sigma^{-1})^o\times(\mathcal{M}\Sigma^{-1})$, with values either in Hot, or in Ind\,Cat\,hot. (\textbf{NB} The two category structures on $\mathcal{M}_0$ deduced from the % [sic] des -> de, « déduites des ces deux foncteurs »
these two functors are canonically isomorphic.)
\[
\begin{tikzcd}[column sep=small]
 & \mathrm{Hot} & (a,b)\longmapsto\mathrm{Ch}^{\Sigma}_{\infty}(a,b)\\
(\mathcal{M}\Sigma^{-1})^o\times(\mathcal{M}\Sigma^{-1}) \arrow[ur] \arrow[dr] & & \\
 & \mathrm{Ind\,Cat\,hot} & (a,b)\longmapsto\mathrm{CH}^{\Sigma}_{\infty}(a,b)
\end{tikzcd}
\]

\msp{page 9}
As the canonical functors
\[
\begin{tikzcd}
\mathrm{Ind\,Cat\,hot} \arrow[r, "\mathrm{Ind}\,\varphi"] & \mathrm{Ind\,Hot} \arrow[r, "\mathrm{Ind}\,\psi"] & \mathrm{Ind\,Hot}_0^{\wedge} \arrow[r, "\varinjlim"] & \mathrm{Hot}_0^{\wedge},\\
\mathrm{Cat\,hot} \arrow[u, hook] \arrow[r, "\varphi"] & \mathrm{Hot} \arrow[u, hook] \arrow[r, "\psi"] & \mathrm{Hot}_0^{\wedge} \arrow[u, hook] \arrow[ur, "\mathrm{id}"'] &
\end{tikzcd}
\]
commuting with finite products, any theory of ``external Hom'' with values in one of these categories, defines one with values in each of those into which it maps.
\printedpage{7}
% head: A. Grothendieck : Dérivateurs, Ch. IX, Édition des Universités de Montpellier II et Paris VII
% foot: Version 12 mars 2002
Besides, these functors are not all fully faithful (this is in any case false for $\varphi$ and $\mathrm{Ind}\,\varphi$, and I do not know what the situation is for $\psi$ and $\mathrm{Ind}\,\psi$ -- the vertical arrows on the other hand are fully faithful (\footnote[3]{and also, $\mathrm{Ind\,Hot}_0^{\wedge}\xrightarrow{\varinjlim}\mathrm{Hot}_0^{\wedge}$ is \textit{not} fully faithful.}).) But they induce bijections for the $\pi_0$. So the category $\MOP$ deduced from a theory of external Hom with values in one of these categories, and the one % [sic] catégoires -> catégories, faute d'orthographe
deduced in another into which it maps, is the same up to canonical isomorphism. As far as this definition of $\MOP$ is concerned, nothing is changed then when we regard the $\mathrm{CH}_{\infty}$ as objects of Ind\,Hot, or only of $\mathrm{Ind\,Hot}_0^{\wedge}$, or even only of $\HotO$, or when we regard the $\mathrm{Ch}_{\infty}$ as objects
\msp{page 10}
in $\HotO$. But of course, the datum of the composition pairings loses precision each time (except perhaps when passing from Hot to $\HotO$, if $\psi$ is fully faithful \dots\,). When we take the images of the $\mathrm{CH}^{\Sigma}_{\infty}(a,b)$ and $\mathrm{Ch}^{\Sigma}_{\infty}(a,b)$ in $\HotO$ (the last category of the diagram), we find canonically isomorphic objects (functorially in $a$, $b$ in $\MOP\simeq\mathcal{M}\Sigma^{-1}$). When we regard them as being respectively in Ind\,Hot and in $\mathrm{Hot}\hookrightarrow\mathrm{Ind\,Hot}$, we find a functorial homomorphism
\[
\mathrm{CH}^{\Sigma}_{\infty}(a,b)\longrightarrow\mathrm{Ch}^{\Sigma}_{\infty}(a,b) \qquad\quad \text{in Ind\,Hot or in } \mathrm{Ind\,Hot}_0^{\wedge} \text{ as we choose.} \tag*{$(*)$}
\]
The $\varinjlim$ of $\mathrm{CH}^{\Sigma}_{\infty}(a,b)$ in Hot does not exist in general (unless the ind-object is essentially constant), but it exists in $\HotO$, and the canonical isomorphism
\[
\varinjlim_{\HotO}\mathrm{CH}^{\Sigma}_{\infty}(a,b)\xrightarrow{\ \sim\ }\mathrm{Ch}^{\Sigma}_{\infty}(a,b) \qquad\quad \text{in } \HotO \tag*{$(**)$}
\]
is the arrow deduced from $(*)$ in $\mathrm{Ind\,Hot}_0^{\wedge}$.

\msp{page 11}
In the two cases CH and Ch (which we are going to denote by the common symbol $C_{\infty}(a,b)$), the functoriality of this object $C_{\infty}(a,b)$ (even in the weakest sense, as an object of $\HotO$) with respect to $a$ and $b$ (even in the weakest sense, as a functor $\mathcal{M}^o\times\mathcal{M}\longrightarrow\HotO$ and not even $(\mathcal{M}\Sigma^{-1})^o\times(\mathcal{M}\Sigma^{-1})\longrightarrow\HotO$) is not at all evident from the direct definition in terms of the explicit direct system in Cat
\[
C_{\infty}(a,b) \;=\; (\mathrm{Ch}_{\tau_0^{2n}}(a,b))_{n\in\mathbf{N}},
\]
for an arrow $a'\longrightarrow a$ or $b\longrightarrow b'$ in $\mathcal{M}$, and even were it in $\Sigma$, in no way defines an arrow in Cat
\[
\begin{array}{rcl}
C_{\infty}(a,b) & \longrightarrow & C_{\infty}(a',b)\\
\text{or}\ \ C_{\infty}(a,b) & \longrightarrow & C_{\infty}(a,b').
\end{array}
\]
This difficulty recalls the one that \textsc{Quillen} had to overcome, in Homotopical Algebra, to face the difference between ``left homotopies'' and ``right homotopies''. I recall how
\printedpage{8}
% head: A. Grothendieck : Dérivateurs, Ch. IX, Édition des Universités de Montpellier II et Paris VII
% foot: Version 12 mars 2002
\msp{page 12}
we can nevertheless define such arrows, as resulting from a (much richer) arrow of ``composition'' of paths between $a'$ and $a$, and $a$ and $b$, or between $a$ and $b$, and $b$ and $b'$, that is in general then the pairing
\[
C_{\infty}(a,b)\times C_{\infty}(b,c)\longrightarrow C_{\infty}(a,c),
\]
whose direct definition is no less problematic, and can in any case \textit{not} be given as an arrow in Cat itself, or in Ind\,Cat. Interpreting e.g. $\mathrm{Ch}_{\infty}(a,b)$ in terms of \textit{infinite} paths (on both sides) from $a$ to $b$, in an obvious sense (but % [sic] cotés -> côtés, accent omis
``essentially finite'' on both sides), we can also define % [sic] cotés -> côtés, accent omis
\[
\mathrm{Ch}^{-}_{\infty}(a,b), \qquad \mathrm{Ch}^{+}_{\infty}(a,b)
\]
as the full subcategories formed respectively of paths which are totally degenerate in degrees $\geq 0$ (i.e. $a_0\longleftarrow a_1\longrightarrow a_2\longleftarrow a_3\cdots$ are all identities), which amounts to paths infinite only to the left, i.e. stopping at $a_0$, and symmetrically for $\mathrm{Ch}^{+}_{\infty}(a,b)$. Same story for the $\mathrm{CH}^{+}$, $\mathrm{CH}^{-}$.

\msp{page 13}
That said, we have this (returning to the ambivalent notion $C_{\infty}$):

\begin{enumerate}
\item[a)] the fully faithful inclusions
\[
C^{-}_{\infty}(a,b)\hookrightarrow C_{\infty}(a,b)\hookleftarrow C^{+}_{\infty}(a,b)
\]
are isomorphisms in Ind\,Cat\,hot resp. in Hot, and
\item[b)] We really do have a composition of paths
\[
C^{-}_{\infty}(a,b)\times C^{+}_{\infty}(b,c)\longrightarrow C_{\infty}(a,c).
\]
\end{enumerate}

The \textit{associativity} in the CH case moreover offers no difficulty, because of the possibility of reducing to the
\[
C_m(a,b)\times C_n(b,c)\times C_p(c,d)\rightrightarrows C_N(a,d)
\]
for $N$ large, and proving equality as arrows in Cat\,hot -- which will follow from the homotopy lemma % [sic] egalité -> égalité, accent omis
on intervals in Cat. In the case of $\mathrm{Ch}_{\infty}$ on the other hand I have not found a proof. (Only for lack of having tried seriously to get down to it? I don't much feel like it, seeing that it might be that $\mathrm{Hot}\longrightarrow\HotO$ is fully faithful, in which case the work becomes useless \dots\,)